Three hundred problems where the graduate toolbox meets numbers, shapes and chance. How well can a typical real number be
approximated by fractions, and how large, in the sense of Hausdorff, is the set of numbers that do better? Why can two
thieves always split a necklace fairly with two cuts? Why does a closed curve in space turn through at least $2\pi$, and
at least $4\pi$ if it is knotted? In how many ways is a cycle a product of transpositions? Each problem mixes ideas from
measure theory, functional and harmonic analysis, topology, Lie groups and algebra, the second part of \emph{Order and
Chance}, and none says which tool to use: finding it is the game. Hints point the way when the search stalls, and every
solution begins with how one might discover the approach before giving a complete proof. Long \emph{Journeys} walk step by
step to famous theorems, from Khinchin's theorem to Kneser's conjecture, and remarks show where each problem leads in
current research.
